\begin{mistake}{2026-04-20}{01}

\begin{problemcontent}
设 \(A = [\alpha_1, \alpha_2, \alpha_3], B = [\beta_1, \beta_2, \beta_3]\) 均为 3 阶矩阵，则下列选项不正确的是（ ）

(A) 若向量组 \(\alpha_1, \alpha_2, \alpha_3\) 与向量组 \(\beta_1, \beta_2, \beta_3\) 等价，则矩阵 \(A\) 与矩阵 \(B\) 等价.

(B) 若 \(r(A) = r(B) = 3\)，则向量组 \(\alpha_1, \alpha_2, \alpha_3\) 与向量组 \(\beta_1, \beta_2, \beta_3\) 等价.

(C) 若 \(r(A) = r(B) = 2\)，则向量组 \(\alpha_1, \alpha_2, \alpha_3\) 与向量组 \(\beta_1, \beta_2, \beta_3\) 等价.

(D) 若 \(r(A, B) = r(B)\)，则向量组 \(\alpha_1, \alpha_2, \alpha_3, \beta_1, \beta_2, \beta_3\) 与向量组 \(\beta_1, \beta_2, \beta_3\) 等价.
\end{problemcontent}

\begin{solutioncontent}
本题选 (C)。

(A) 正确。向量组等价的必要条件是它们的秩相等，即 \(r(A) = r(B)\)。同型矩阵秩相等即为矩阵等价。

(B) 正确。\(r(A) = r(B) = 3\) 且矩阵为 3 阶，说明两向量组均构成 3 维空间 \(\mathbb{R}^3\) 的基。它们张成的空间都是整个 \(\mathbb{R}^3\)，必然完全重合，故能互相表出，向量组等价。

(C) 错误。\(r(A) = r(B) = 2\) 仅说明两个向量组张成的空间维数均为 2，但不能保证空间完全重合（例如 \(\mathbb{R}^3\) 中两个不同且相交的平面），里面的向量不一定能相互表出，故不一定等价。

(D) 正确。\(r(A, B) = r(B)\) 等价于向量组 \(\alpha_i\) 能由 \(\beta_i\) 线性表出。因此，合并后的“大组”自然能全部由“小组” \(\beta_i\) 表出；反之，“小组”作为“大组”的子集，天然能由“大组”表出。二者互相表出，故等价。
\end{solutioncontent}

\mistakeknowledge{
\begin{itemize}
 \item \textbf{核心定理}：向量组等价 \(\iff\) 相互表出 \(\iff r(A)=r(B)=r(A,B)\)；同型矩阵等价 \(\iff r(A)=r(B)\)。
 \item \textbf{易错点}：误将 \(r(A)=r(B)\) 等同于向量组等价。秩相等仅代表生成空间的维数相同，只有满秩或满足 \(r(A,B)=r(B)\) 时才能证明空间完全重合。
 \item \textbf{关联知识}：\(r(A,B)=r(B)\) 的本质含义是矩阵 \(A\) 的列向量能被矩阵 \(B\) 的列向量线性表出；任何子向量组均可由包含它的全集向量组线性表出。
\end{itemize}}

\end{mistake}
